\section{The Spinglass model}
We now move on a more complex model, that is the Edwards–Anderson model.
\subsection{The Edwards–Anderson model}
The Edwards–Anderson (EA) model -- first introduced by .. \cite{SFEdwards_1975} -- is a finite dimensional spin glass model with short range interaction. It is not a mean field spin glass model with infinite range interactions such as the Sherrington–Kirkpatrick model \cite{SK_model}. The model is formally described as follows:
\begin{equation}
  H_\text{EA}=\sum_{\langle i,j \rangle}J_{ij}S_iS_j.
  \label{eq:EA_model}
\end{equation}
Where $J_{ij}$ is a random variable taking its values in $\{-1,\,1\}$ with equal probability. Thus, it can easily be represented by a \Bits.
Under a spin flip ($S_k \to - S_k$), the energy shift is:
\begin{equation}
  \Delta E_k=2 S_k \sum_{i \partial k}J_{ki}S_i.
  \label{eq:EA_energy_shift}
\end{equation}
Performing the same derivation as for the Ising model, the spin flip condition for this model is thus:
\begin{equation}
  \lfloor \rho \rfloor \geq S_k  \sum_{i \in \partial k}J_{ki}S_i,\\
  \label{eq:flip_EA_int}
\end{equation}
with $\rho$ the same exponential random variable as in the previous section.
\subsection{Bitwise Implementation}
\subsection{Performance test}
